%! Setting Up the SVD — Unit 7, Lesson 7.0 — Homework (Answer Key)
\documentclass[10pt]{article}
\usepackage{linalg-article}
\usepackage{linalg-key}   % pulls in -boxes; do NOT also load -boxes
\newcommand{\TallMath}[1]{$\displaystyle #1\rule[-1.4em]{0pt}{3.2em}$}
\newcommand{\uu}{\mathbf{u}}
\newcommand{\vv}{\mathbf{v}}
\newcommand{\ee}{\mathbf{e}}
\newcommand{\zero}{\mathbf{0}}
\newcommand{\T}{^{\top}}

\begin{document}

\pageheader{Unit 7, Lesson 7.0}{Homework --- Answer Key}
\namedateperiod

\vspace{0.12in}

\begin{notesbox}{Practice}
A matrix stretches the unit circle into an \textbf{ellipse}. The \textbf{singular values}
$\sigma_1\ge\sigma_2\ge0$ are its half-widths; a perpendicular input pair $\vv_1,\vv_2$ maps to a
perpendicular output pair with $A\vv_i=\sigma_i\uu_i$.

\begin{enumerate}[leftmargin=1.9em, itemsep=12pt]
  \item \textbf{Read $\sigma$ off a diagonal matrix.} Give the singular values $\sigma_1,\sigma_2$
        (larger first) and describe the ellipse the unit circle becomes.
    \begin{enumerate}[label=(\alph*), leftmargin=1.8em, itemsep=4pt]
      \item $\begin{bmatrix} 5 & 0 \\ 0 & 2 \end{bmatrix}$ \hfill \ans{$\sigma_1=5,\ \sigma_2=2$; ellipse with half-widths $5$ ($x$) and $2$ ($y$).}
      \item $\begin{bmatrix} 1 & 0 \\ 0 & 6 \end{bmatrix}$ \hfill \ans{$\sigma_1=6,\ \sigma_2=1$; half-widths $6$ ($y$) and $1$ ($x$).}
      \item $\begin{bmatrix} 3 & 0 \\ 0 & 3 \end{bmatrix}$ \hfill \ans{$\sigma_1=3,\ \sigma_2=3$; the circle just scales to a circle of radius $3$.}
    \end{enumerate}
  \item \textbf{Perpendicular in, perpendicular out.} For $S=\begin{bmatrix} 0 & 4 \\ 3 & 0 \end{bmatrix}$,
        compute $S\ee_1$ and $S\ee_2$. Show the outputs are perpendicular, give $\sigma_1,\sigma_2$, and
        state the output direction $\uu_1$ belonging to $\sigma_1$. Is $\uu_1$ the same as its input
        direction?
        \ansline{$S\ee_1=\begin{bsmallmatrix} 0\\3 \end{bsmallmatrix}$ (length $3$), $S\ee_2=\begin{bsmallmatrix} 4\\0 \end{bsmallmatrix}$ (length $4$); dot $=0$ so perpendicular. $\sigma_1=4$, $\sigma_2=2$? \textbf{No:} $\sigma_1=4,\ \sigma_2=3$. The largest stretch ($4$) comes from input $\vv_1=\ee_2$, output $\uu_1=\begin{bsmallmatrix} 1\\0 \end{bsmallmatrix}$ --- \emph{not} the same as $\vv_1=\ee_2$.}
  \item \textbf{A symmetric matrix (recall Unit~6).} For $A=\begin{bmatrix} 3 & 1 \\ 1 & 3 \end{bmatrix}$,
        the vectors $\begin{bmatrix} 1 \\ 1 \end{bmatrix}$ and $\begin{bmatrix} 1 \\ -1 \end{bmatrix}$ are
        perpendicular eigenvectors. Compute $A\begin{bmatrix} 1 \\ 1 \end{bmatrix}$ and
        $A\begin{bmatrix} 1 \\ -1 \end{bmatrix}$, and read off the singular values.
        \ansline{$A\begin{bsmallmatrix} 1\\1 \end{bsmallmatrix}=\begin{bsmallmatrix} 4\\4 \end{bsmallmatrix}=4\begin{bsmallmatrix} 1\\1 \end{bsmallmatrix}$ and $A\begin{bsmallmatrix} 1\\-1 \end{bsmallmatrix}=\begin{bsmallmatrix} 2\\-2 \end{bsmallmatrix}=2\begin{bsmallmatrix} 1\\-1 \end{bsmallmatrix}$, so $\sigma_1=4$, $\sigma_2=2$.}
  \item \textbf{Why never negative?} Explain why a singular value can never be negative, and why the
        singular values are listed from largest to smallest.
        \ansline{A singular value is a \emph{length} --- a half-width of the ellipse / a stretch factor --- and lengths are never negative. Listing largest first names $\sigma_1$ as the biggest stretch (the long axis), so the ordering is unambiguous.}
  \item \textbf{Why the SVD wins.} In a sentence or two, describe what $A=U\Sigma V\T$ says a matrix does
        (rotate, stretch, rotate), and state one reason the SVD is more widely usable than eigenvectors.
        \ansline{$A=U\Sigma V\T$: first rotate to the input singular frame ($V\T$), then stretch each axis by a singular value ($\Sigma$), then rotate to the output frame ($U$). Unlike eigenvectors, \emph{every} matrix has an SVD --- including non-square ones --- and the singular values are always real and $\ge0$.}
\end{enumerate}
\end{notesbox}

\begin{extensionbox}
\textbf{Finding singular values with $A\T A$ (a first look at Lesson~7.1).} The singular values are
$\sigma=\sqrt{\lambda}$, where $\lambda$ are the eigenvalues of the symmetric matrix $A\T A$. Use
$S=\begin{bmatrix} 0 & 4 \\ 3 & 0 \end{bmatrix}$ from Problem~2.
\begin{enumerate}[label=(\alph*), leftmargin=1.8em, itemsep=5pt]
  \item Compute $A\T A$ (remember $A\T$ swaps rows and columns).
  \item Find the eigenvalues $\lambda$ of $A\T A$, then take $\sigma=\sqrt{\lambda}$. Do they match the
        $\sigma$'s you read off in Problem~2?
\end{enumerate}
\ansline{(a) $A\T=\begin{bsmallmatrix} 0 & 3\\ 4 & 0 \end{bsmallmatrix}$, so $A\T A=\begin{bsmallmatrix} 9 & 0\\ 0 & 16 \end{bsmallmatrix}$.}
\ansline{(b) Eigenvalues $\lambda=16$ and $\lambda=9$ (diagonal), so $\sigma=\sqrt{16}=4$ and $\sqrt{9}=3$ --- matching $\sigma_1=4,\ \sigma_2=3$ from Problem~2.}
\end{extensionbox}

\begin{spiralbox}
\textbf{Coming next --- Lesson 7.1: Singular Values and Singular Vectors.} Today you \emph{read} singular
values off nice matrices. Next we \emph{find} them for any matrix: the eigenvalues of $A\T A$ give
$\sigma^2$, and its eigenvectors give the input directions $\vv_i$ --- then $\uu_i=A\vv_i/\sigma_i$.
\end{spiralbox}

\begin{teachernote}
Item~1 drills reading $\sigma$ off diagonals, including the ordering convention (part~b: $\sigma_1=6$
comes from the \emph{second} axis) and the ``circle stays a circle'' case (part~c). Item~2 is the
input-$\ne$-output-frame idea on a swap matrix (note the deliberate ``No'' correcting a wrong $\sigma_2=2$
to $3$). Item~3 ties singular values to Unit~6 (perpendicular eigenvectors of a symmetric matrix are the
singular vectors; positive eigenvalues are the singular values). Items~4--5 are the target justifications
($\sigma\ge0$ as a length; SVD exists for every matrix). The extension is the Lesson~7.1 method
($\sigma=\sqrt{\lambda}$ from $A\T A$). Most common slips: forgetting to order $\sigma_1\ge\sigma_2$;
conflating input $\vv_i$ and output $\uu_i$; reporting a length as negative.
\end{teachernote}

\end{document}
